\section{Verification Methodology}
\label{sec:validation}

Because the central premise of the mod is that its behavior is not a scripted
approximation, the simulator was checked during development against nine closed-form and
statistical properties, each executable as a plain Java program requiring no game instance
(the \texttt{sim} package has no dependency on any Minecraft class at all, exactly as
CircuitCraft's solver package does~\cite{picos2026circuitcraft}): that $X|0\rangle$ produces
exactly $|1\rangle$; that a Hadamard applied to $|0\rangle$ produces exactly a 50/50
superposition; that the two-qubit Bell-state and three-qubit GHZ-state circuits of
Section~\ref{sec:experiments} produce their exact expected amplitudes; that a Toffoli gate
reproduces the reversible-AND truth table for all eight classical three-bit inputs; that the
phase-kickback circuit of Section~\ref{sec:exp-kickback} produces a deterministic $|1\rangle$;
that 20{,}000 simulated measurements of a Bell state never once produced a $|01\rangle$ or
$|10\rangle$ outcome and converged to within half a percentage point of the predicted 50\%
split; and that total probability remains exactly $1$ after a long sequence of mixed single-
and multi-qubit gates, confirming that every gate application preserves unitarity. All nine
checks passed.

Because the mod targets a version of Minecraft (26.2) released after the training data
cutoff of the large language model used to assist in its development, and because that
version introduced several breaking changes to modding-relevant application programming
interfaces, every such interface used by the \texttt{block}, \texttt{blockentity}, and
\texttt{network} packages (block registration, the two-phase heads-up-display rendering
interface, and the per-item model-definition file format) was checked directly against the
decompiled game source for this specific version prior to being relied upon, following the
same methodological practice established during CircuitCraft's own
development~\cite{picos2026circuitcraft}: verifying an assumption against the authoritative
source for the exact version in use, rather than against generic recollection of how
equivalent interfaces behaved in earlier versions of the game.

For transparency, it should be noted that, at the time of writing, the nine checks described
above were performed ad hoc during development rather than captured as a persisted,
automatically re-run test suite within the repository -- the same caveat, and the same
remedy, noted in CircuitCraft's own validation section~\cite{picos2026circuitcraft}.
Formalizing these checks into a checked-in regression suite is planned prior to a longer-term
evaluation of the tool (Section~\ref{sec:limitations}) and would strengthen a reviewer's
ability to independently confirm the accuracy claims made here.
